% GENERATED by build_do_now_tex.py on 2026-04-24 --- do not hand-edit; regenerate instead.
\documentclass[11pt]{article}
\usepackage{preamble}

\begin{document}

\packetheading{Lesson 4-4 \textperiodcentered\ Period 1 \textperiodcentered\ Do Now}{L44 P1}

\vspace{0.25em}
\noindent\textbf{Name:}\ \rule{2in}{0.4pt}\quad
\textbf{Period:}\ \rule{0.5in}{0.4pt}\quad
\textbf{Date:}\ \rule{1in}{0.4pt}
\vspace{0.5em}

\bankitem{Do Now}{%
CRITIQUE & EXPLAIN
Teo and Evander find the following exercise in their homework:
 (1)/(2) + (1)/(3) + (1)/(9) 

[IMAGE: A tablet displaying the math problem 1/2 + 1/3 + 1/9]

\begin{enumerate}
 \item[A.] Teo claims that a common denominator of the sum is 2 + 3 + 9 = 14. Evander claims that it is 2 · 3 · 9 = 54. Is either student correct? Explain why or why not.
 \item[B.] Find the sum, explaining the method you use.
 \item[C.] Construct Arguments Timothy states that the quickest way to find the sum of any two fractions with unlike denominators is to multiply their denominators to find a common denominator, and then rewrite each fraction with that denominator. Do you agree?
\end{enumerate}
}{}
\vspace{2.5in}

\end{document}
